% page 349 of the facsimile (image p-348 of the scan)
% block 160.0 x 242.0 mm ; left margin 42.0 mm
\pgc{20.320mm}{0.000mm}{
\l{\cel{54}341}
\l{}
\l{}
\l{\sou{listo}.\cel{2}Serio de nomi di personi o di kozi, de numeri, de}
\l{\cel{3}nombri, e c., skribita ica dop ita. - DEFIRS.}
\l{}
\l{listelo. (\sou{arkitekt.}) Mikra muluro quadrata e plata, qua se-}
\l{\cel{3}paras du altri plu granda e plu salianta - ed uzesas por}
\l{\cel{3}reprezentar kanalo sur pilastro, o kolono. - DEFIRS.}
\l{}
\l{\sou{lito}. I. Moblo destinita a recevor su-kushonti. - II. Parto}
\l{\cel{3}kava, plu o min profunda, en qua kontenesas e fluas la aquo}
\l{\cel{3}di fluvio, di rivero, di revereto. - FI.}
\l{}
\l{\sou{litanio}. I. (\sou{religio katol.)}\cel{2}Prego forme di serio de advoki}
\l{\cel{3}a Iesu Kristo, a la santa Virgino, e c., en qui la sama for-}
\l{\cel{3}mulo di supliko repetesas multafoye. (Uzesas ordinare en}
\l{\cel{3}pluralo.) - II.Enumero tedoza. - DEFIRS.}
\l{}
\l{\sou{litargiro.(kemio}).Protoxido di plombo fuzita, mi-vitrigita,}
\l{\cel{3}uzata por preparar la plombo-sali qui esas un la kompozanti}
\l{\cel{3}di kristalo e per qua onu falsigis ula vino-sorti por dimi-}
\l{\cel{3}nutar lia grado di acideso. - DEFIRS.}
\l{}
\l{\sou{litero}. Signo alfabetala per qua onu reprezentas singla son-}
\l{\cel{3}uno di linguo. - DEFIRS.}
\l{}
\l{\sou{literaturo}. Ensemblo di la stilo-verki (en qui preponderas}
\l{\cel{3}la beleso di kompozo, di expreso : poezio, eloquenteso,}
\l{\cel{3}e c. ) di yarcento, di naciono. - DEFIRS.}
\l{}
\l{\sou{litio}. (\sou{kemio}).Korpo simpla, la maxim lejera ek omna metali}
\l{\cel{3}konocata. -\cel{1}\sou{Simb. kemiala}\cel{1}: Li. - DEFIRS.}
\l{}
\l{litiero. - Lito ek palio, folii, filiko sika, por la animali,}
\l{\cel{3}en la stabli. - EFI.}
\l{}
\l{litijo. (\sou{yurocienco})Kontesto-punto qua divenas la materio di}
\l{\cel{3}proceso. -DEFIS.}
\l{}
\l{litografar. (\sou{trans}.)\cel{2}Reproduktar per preso la skriburi, la}
\l{\cel{3}desegnuri, per trasar li (per krayono od inko grasoza) sur}
\l{\cel{3}petro kalkoza di qua la grano esas tre mikra. - DEFIS.}
\l{}
\l{litologio. Ta parto di la historio naturala,qua koncernas la}
\l{\cel{3}petri. - DEFIRS.}
\l{}
\l{litoro. La regiono di qua ula parto di lua konturo esas an}
\l{\cel{3}maro. - DEFIS.}
\l{}
\l{litoto. (\sou{retoriko}).Figuro per qua onu minfortigas la expreso}
\l{\cel{3}di sua penso, por komprenigar la maxim multo per la minim}
\l{\cel{3}multo. - DEFIS.}
\l{}
\l{litotomiar. (trans.\cel{2}(kirurg.) Dividar ed extraktar (per}
\l{\cel{3}operaco) kalkolo\cel{2}ek la veziko. -\cel{2}DEFIRS.}
}
